%============================================================ 영문 표제·초록
%  EAAI 는 영문 저널이다. 이 한글 원고는 투고본의 기준 원고이므로, 번역 단계에서
%  용어와 수치가 흔들리지 않도록 영문 제목·초록·주제어를 여기에 고정해 둔다.
%  ★ 수치는 국문 초록과 한 글자도 다르면 안 된다 — tools/check_paper_meta.py 가 대조한다.
%  ★ 이 파일은 투고 원고에 들어가지 않는다 — elsarticle 에는 본문 뒤 제2 초록에
%    대응하는 요소가 없다. 독립 문서 영문초록_translation_reference.tex 가 읽는다.

\section*{English Abstract}
\addcontentsline{toc}{section}{English Abstract}

% 14.4pt 볼드 한 줄로는 판면을 넘는다. 줄바꿈을 손으로 준다.
\noindent\textbf{\large LLM-Based Command Architecture and U-Net Score-Map\\
Multi-Agent Cooperative Reinforcement Learning\\
for Defense Against Suicide USV Swarms}

\vspace{4pt}
\noindent
Low-cost suicide unmanned surface vehicle (USV) swarms evade conventional counters:
electronic warfare cannot reach wire-guided targets, and interceptors cost more than their
targets and saturate. We propose a two-layer defense in which a few boats slower than the
attackers pre-block the mothership approach corridor with net walls. In the strategy layer a
large language model (LLM) commander assigns each boat to an enemy cluster; in the maneuver
layer a U-Net emits a score map over a rasterized battlefield grid and each boat points at
pixels to place its net wall. Both layers update on the same period but differ sharply in
response latency, so only the slow strategy layer is decoupled asynchronously. Pixel pointing
confines the action space to the grid, imposes safety constraints by masking, and is
invariant to swarm size; training extends critic-free group relative policy optimization to
the multi-agent case. With no combinatorial post-optimization, we collected
720 seed-paired episodes (8 conditions $\times$ 3 formations $\times$ 30 seeds) and separated
the two layers' contributions in a $2\times2$ decomposition. The maneuver layer alone raised
capture rate from $0.844$ to $0.886$; the full system cut nets per capture by 44.5\,\% and total
turning by 46.6\,\% at the same defense level, and the two contributions were \emph{additive}
rather than synergistic. The commander paid off only above a capability threshold --- a cloud
model gained $+0.060$, whereas local Qwen2.5 7B was significantly worse than the code heuristic
at $-0.192$. Commander performance was explained by plan stability across replans, not by
assignment coverage: in periodically replanning architectures, temporal consistency matters
more than local optimality.

\vspace{4pt}
\noindent\textbf{Keywords}\;:\;
unmanned surface vehicle; swarm defense; multi-agent reinforcement learning; task assignment;
large language model; semantic segmentation; group relative policy optimization

\vspace{6pt}
\noindent\textbf{Highlights}
\begin{itemize}[leftmargin=1.1em, itemsep=1pt, topsep=2pt, parsep=0pt]
\item Two judgment layers share one 25 s period; only the slow one is decoupled
\item Pointing at score-map pixels keeps the action space invariant to swarm size
\item Critic-free, heuristic-anchored group relative optimization for multi-boat coordination
\item Across 720 paired episodes the two layers are additive in 9 of 10 metrics
\item Commander skill is explained by plan stability, not assignment coverage
\end{itemize}
